\begin{abstract}
本实验使用AI编程工具制作了一个包含自我介绍、数据可视化和交互彩蛋的个人主页，部署于GitHub Pages。
实验分为两个阶段：第一阶段使用Codex采用"逐步构建"策略（先搭骨架、再加模块、逐模块调整效果），在搭建骨架和加入开屏动画后即因缺乏核心设计理念和对设计方向的掌控感而停滞；第二阶段改用Claude Code采用"设计优先"策略——先与AI深度讨论完整的设计计划（调性、色彩、字体、六大区域、动效规格），汇总为设计文档，再按文档执行代码实现。最终打造了以"温润材质 $\times$ 锐利工艺 $\times$ 生命律动"为调性的完整主页，包含四个主题数据可视化面板、横向翻阅招牌交互、三重页脚彩蛋和隐藏贪吃蛇游戏。本实验通过两种策略的对比，论证了Vibe Coding中"先设计、后编码"策略对非技术背景用户的优越性。
\end{abstract}
